% Required in the manuscript preamble: \usepackage{booktabs}
% Include with: % Required in the manuscript preamble: \usepackage{booktabs}
% Include with: % Required in the manuscript preamble: \usepackage{booktabs}
% Include with: % Required in the manuscript preamble: \usepackage{booktabs}
% Include with: \input{tables/table1.tex}
% Source definitions and cohort/eligibility differences: tables/README.md.
\begin{table*}[t]
\centering
\caption{Prediction accuracy and computational readiness. Part A uses
raw model feasibility labels; Part B reports archived warmed computation
times and realized prediction availability. Cohorts and missing
measurements are distinguished below.}
\label{tab:prediction_accuracy_readiness}
\begingroup
\scriptsize
\setlength{\tabcolsep}{3.5pt}
\renewcommand{\arraystretch}{1.15}
\textbf{A. Feasibility prediction and sampled first-loss errors}\par\smallskip
\begin{tabular}{@{}lrrrrrrrr@{}}
\toprule
Mission / predicted property &
\shortstack{Scenarios\\$N$} &
\shortstack{Label\\pairs} &
\shortstack{Accuracy\\(\%)} &
\shortstack{False-feasible\\$n$ (\%)} &
\shortstack{False-reject\\$n$ (\%)} &
\shortstack{Boundary\\MAE (ms)} &
\shortstack{Boundary\\median AE (ms)} &
\shortstack{Boundary\\pairs} \\
\midrule
Inspection Descent: command & 19 & 209 & 98.09 & 2 (1.56) & 2 (2.47) & 16.67 & 0.00 & 18 \\
Geofence Turn: command & 20 & 380 & 99.21 & 1 (0.63) & 2 (0.90) & 3.16 & 0.00 & 19 \\
\midrule
Underpass: command & 400 & 7600 & 99.29 & 25 (1.28) & 29 (0.51) & 11.45 & 0.00 & 318 \\
Underpass: safe completion & 400 & 7600 & 99.29 & 25 (1.28) & 29 (0.51) & 11.45 & 0.00 & 318 \\
Underpass: recovery & 400 & 7600 & 99.61 & 12 (2.18) & 18 (0.26) & 7.17 & 0.00 & 237 \\
\bottomrule
\end{tabular}

\medskip
\textbf{B. Evaluation-cohort timing and readiness}\par\smallskip
\begin{tabular}{@{}lrrrrr@{}}
\toprule
Mission / predictor &
\shortstack{Timed\\forecasts} &
\shortstack{Runtime\\median (ms)} &
\shortstack{Runtime\\p95 (ms)} &
\shortstack{Ready by scheduled dispatch\\$n/N$ (\%)} &
\shortstack{Late relative to physical\\recovery loss ($n/N$)} \\
\midrule
Inspection Descent: command-only & 276 & 42.26 & 61.39 & 2760/3036 (90.91) & N/M \\
Geofence Turn: command-only & 300 & 83.70 & 104.29 & 5480/6900 (79.42) & N/M \\
\midrule
Underpass: command-only & 400 & 142.36 & 146.71 & 3600/5200 (69.23) & 0/243 \\
Underpass: RecoveryAware & 400 & 149.12 & 166.96 & 3537/5200 (68.02) & 0/243 \\
Underpass: \textbf{full} & 400 & 149.00 & 167.44 & 3547/5200 (68.21) & 0/243 \\
\bottomrule
\end{tabular}

\smallskip
\begin{minipage}{\textwidth}
\scriptsize
\textit{Definitions and scope.}
False-feasible and false-reject rates divide by physically infeasible and
feasible label counts, respectively. AE denotes absolute error.
Part A uses 19 Inspection and 20 Geofence \emph{development} scenarios,
and 400 Underpass \emph{evaluation} scenarios; Part B uses 276, 300 and
400 evaluation scenarios. Boundary errors are sampled first-loss errors,
not continuous-time deadline errors. Original eligibility conventions are
retained: Inspection excludes one initially invalid forecast; Underpass
includes one initial failure as sample zero for command/completion error.
Unobserved boundaries receive no finite error. Underpass command/completion
and recovery losses are unobserved in 82 and 157 scenarios, respectively.
N/M means not measured, not zero. Independent completion/recovery prediction
accuracy is not reported for Inspection or Geofence. Late availability
means at/after the first physically unsafe recovery sample, among the 243
Underpass scenarios with observed recovery loss. All 400 Underpass full
forecasts become ready at a tested safe recovery start.
Readiness uses the original control-step scheduling; timing was measured
offline and replayed on the simulator clock. Ready-by-dispatch counts
include cases canceled earlier. Numerical profiles coincide for
RecoveryAware and full StaleGuard, but their runtime measurements are independent.
\end{minipage}
\endgroup
\end{table*}

% Source definitions and cohort/eligibility differences: tables/README.md.
\begin{table*}[t]
\centering
\caption{Prediction accuracy and computational readiness. Part A uses
raw model feasibility labels; Part B reports archived warmed computation
times and realized prediction availability. Cohorts and missing
measurements are distinguished below.}
\label{tab:prediction_accuracy_readiness}
\begingroup
\scriptsize
\setlength{\tabcolsep}{3.5pt}
\renewcommand{\arraystretch}{1.15}
\textbf{A. Feasibility prediction and sampled first-loss errors}\par\smallskip
\begin{tabular}{@{}lrrrrrrrr@{}}
\toprule
Mission / predicted property &
\shortstack{Scenarios\\$N$} &
\shortstack{Label\\pairs} &
\shortstack{Accuracy\\(\%)} &
\shortstack{False-feasible\\$n$ (\%)} &
\shortstack{False-reject\\$n$ (\%)} &
\shortstack{Boundary\\MAE (ms)} &
\shortstack{Boundary\\median AE (ms)} &
\shortstack{Boundary\\pairs} \\
\midrule
Inspection Descent: command & 19 & 209 & 98.09 & 2 (1.56) & 2 (2.47) & 16.67 & 0.00 & 18 \\
Geofence Turn: command & 20 & 380 & 99.21 & 1 (0.63) & 2 (0.90) & 3.16 & 0.00 & 19 \\
\midrule
Underpass: command & 400 & 7600 & 99.29 & 25 (1.28) & 29 (0.51) & 11.45 & 0.00 & 318 \\
Underpass: safe completion & 400 & 7600 & 99.29 & 25 (1.28) & 29 (0.51) & 11.45 & 0.00 & 318 \\
Underpass: recovery & 400 & 7600 & 99.61 & 12 (2.18) & 18 (0.26) & 7.17 & 0.00 & 237 \\
\bottomrule
\end{tabular}

\medskip
\textbf{B. Evaluation-cohort timing and readiness}\par\smallskip
\begin{tabular}{@{}lrrrrr@{}}
\toprule
Mission / predictor &
\shortstack{Timed\\forecasts} &
\shortstack{Runtime\\median (ms)} &
\shortstack{Runtime\\p95 (ms)} &
\shortstack{Ready by scheduled dispatch\\$n/N$ (\%)} &
\shortstack{Late relative to physical\\recovery loss ($n/N$)} \\
\midrule
Inspection Descent: command-only & 276 & 42.26 & 61.39 & 2760/3036 (90.91) & N/M \\
Geofence Turn: command-only & 300 & 83.70 & 104.29 & 5480/6900 (79.42) & N/M \\
\midrule
Underpass: command-only & 400 & 142.36 & 146.71 & 3600/5200 (69.23) & 0/243 \\
Underpass: RecoveryAware & 400 & 149.12 & 166.96 & 3537/5200 (68.02) & 0/243 \\
Underpass: \textbf{full} & 400 & 149.00 & 167.44 & 3547/5200 (68.21) & 0/243 \\
\bottomrule
\end{tabular}

\smallskip
\begin{minipage}{\textwidth}
\scriptsize
\textit{Definitions and scope.}
False-feasible and false-reject rates divide by physically infeasible and
feasible label counts, respectively. AE denotes absolute error.
Part A uses 19 Inspection and 20 Geofence \emph{development} scenarios,
and 400 Underpass \emph{evaluation} scenarios; Part B uses 276, 300 and
400 evaluation scenarios. Boundary errors are sampled first-loss errors,
not continuous-time deadline errors. Original eligibility conventions are
retained: Inspection excludes one initially invalid forecast; Underpass
includes one initial failure as sample zero for command/completion error.
Unobserved boundaries receive no finite error. Underpass command/completion
and recovery losses are unobserved in 82 and 157 scenarios, respectively.
N/M means not measured, not zero. Independent completion/recovery prediction
accuracy is not reported for Inspection or Geofence. Late availability
means at/after the first physically unsafe recovery sample, among the 243
Underpass scenarios with observed recovery loss. All 400 Underpass full
forecasts become ready at a tested safe recovery start.
Readiness uses the original control-step scheduling; timing was measured
offline and replayed on the simulator clock. Ready-by-dispatch counts
include cases canceled earlier. Numerical profiles coincide for
RecoveryAware and full StaleGuard, but their runtime measurements are independent.
\end{minipage}
\endgroup
\end{table*}

% Source definitions and cohort/eligibility differences: tables/README.md.
\begin{table*}[t]
\centering
\caption{Prediction accuracy and computational readiness. Part A uses
raw model feasibility labels; Part B reports archived warmed computation
times and realized prediction availability. Cohorts and missing
measurements are distinguished below.}
\label{tab:prediction_accuracy_readiness}
\begingroup
\scriptsize
\setlength{\tabcolsep}{3.5pt}
\renewcommand{\arraystretch}{1.15}
\textbf{A. Feasibility prediction and sampled first-loss errors}\par\smallskip
\begin{tabular}{@{}lrrrrrrrr@{}}
\toprule
Mission / predicted property &
\shortstack{Scenarios\\$N$} &
\shortstack{Label\\pairs} &
\shortstack{Accuracy\\(\%)} &
\shortstack{False-feasible\\$n$ (\%)} &
\shortstack{False-reject\\$n$ (\%)} &
\shortstack{Boundary\\MAE (ms)} &
\shortstack{Boundary\\median AE (ms)} &
\shortstack{Boundary\\pairs} \\
\midrule
Inspection Descent: command & 19 & 209 & 98.09 & 2 (1.56) & 2 (2.47) & 16.67 & 0.00 & 18 \\
Geofence Turn: command & 20 & 380 & 99.21 & 1 (0.63) & 2 (0.90) & 3.16 & 0.00 & 19 \\
\midrule
Underpass: command & 400 & 7600 & 99.29 & 25 (1.28) & 29 (0.51) & 11.45 & 0.00 & 318 \\
Underpass: safe completion & 400 & 7600 & 99.29 & 25 (1.28) & 29 (0.51) & 11.45 & 0.00 & 318 \\
Underpass: recovery & 400 & 7600 & 99.61 & 12 (2.18) & 18 (0.26) & 7.17 & 0.00 & 237 \\
\bottomrule
\end{tabular}

\medskip
\textbf{B. Evaluation-cohort timing and readiness}\par\smallskip
\begin{tabular}{@{}lrrrrr@{}}
\toprule
Mission / predictor &
\shortstack{Timed\\forecasts} &
\shortstack{Runtime\\median (ms)} &
\shortstack{Runtime\\p95 (ms)} &
\shortstack{Ready by scheduled dispatch\\$n/N$ (\%)} &
\shortstack{Late relative to physical\\recovery loss ($n/N$)} \\
\midrule
Inspection Descent: command-only & 276 & 42.26 & 61.39 & 2760/3036 (90.91) & N/M \\
Geofence Turn: command-only & 300 & 83.70 & 104.29 & 5480/6900 (79.42) & N/M \\
\midrule
Underpass: command-only & 400 & 142.36 & 146.71 & 3600/5200 (69.23) & 0/243 \\
Underpass: RecoveryAware & 400 & 149.12 & 166.96 & 3537/5200 (68.02) & 0/243 \\
Underpass: \textbf{full} & 400 & 149.00 & 167.44 & 3547/5200 (68.21) & 0/243 \\
\bottomrule
\end{tabular}

\smallskip
\begin{minipage}{\textwidth}
\scriptsize
\textit{Definitions and scope.}
False-feasible and false-reject rates divide by physically infeasible and
feasible label counts, respectively. AE denotes absolute error.
Part A uses 19 Inspection and 20 Geofence \emph{development} scenarios,
and 400 Underpass \emph{evaluation} scenarios; Part B uses 276, 300 and
400 evaluation scenarios. Boundary errors are sampled first-loss errors,
not continuous-time deadline errors. Original eligibility conventions are
retained: Inspection excludes one initially invalid forecast; Underpass
includes one initial failure as sample zero for command/completion error.
Unobserved boundaries receive no finite error. Underpass command/completion
and recovery losses are unobserved in 82 and 157 scenarios, respectively.
N/M means not measured, not zero. Independent completion/recovery prediction
accuracy is not reported for Inspection or Geofence. Late availability
means at/after the first physically unsafe recovery sample, among the 243
Underpass scenarios with observed recovery loss. All 400 Underpass full
forecasts become ready at a tested safe recovery start.
Readiness uses the original control-step scheduling; timing was measured
offline and replayed on the simulator clock. Ready-by-dispatch counts
include cases canceled earlier. Numerical profiles coincide for
RecoveryAware and full StaleGuard, but their runtime measurements are independent.
\end{minipage}
\endgroup
\end{table*}

% Source definitions and cohort/eligibility differences: tables/README.md.
\begin{table*}[t]
\centering
\caption{Prediction accuracy and computational readiness. Part A uses
raw model feasibility labels; Part B reports archived warmed computation
times and realized prediction availability. Cohorts and missing
measurements are distinguished below.}
\label{tab:prediction_accuracy_readiness}
\begingroup
\scriptsize
\setlength{\tabcolsep}{3.5pt}
\renewcommand{\arraystretch}{1.15}
\textbf{A. Feasibility prediction and sampled first-loss errors}\par\smallskip
\begin{tabular}{@{}lrrrrrrrr@{}}
\toprule
Mission / predicted property &
\shortstack{Scenarios\\$N$} &
\shortstack{Label\\pairs} &
\shortstack{Accuracy\\(\%)} &
\shortstack{False-feasible\\$n$ (\%)} &
\shortstack{False-reject\\$n$ (\%)} &
\shortstack{Boundary\\MAE (ms)} &
\shortstack{Boundary\\median AE (ms)} &
\shortstack{Boundary\\pairs} \\
\midrule
Inspection Descent: command & 19 & 209 & 98.09 & 2 (1.56) & 2 (2.47) & 16.67 & 0.00 & 18 \\
Geofence Turn: command & 20 & 380 & 99.21 & 1 (0.63) & 2 (0.90) & 3.16 & 0.00 & 19 \\
\midrule
Underpass: command & 400 & 7600 & 99.29 & 25 (1.28) & 29 (0.51) & 11.45 & 0.00 & 318 \\
Underpass: safe completion & 400 & 7600 & 99.29 & 25 (1.28) & 29 (0.51) & 11.45 & 0.00 & 318 \\
Underpass: recovery & 400 & 7600 & 99.61 & 12 (2.18) & 18 (0.26) & 7.17 & 0.00 & 237 \\
\bottomrule
\end{tabular}

\medskip
\textbf{B. Evaluation-cohort timing and readiness}\par\smallskip
\begin{tabular}{@{}lrrrrr@{}}
\toprule
Mission / predictor &
\shortstack{Timed\\forecasts} &
\shortstack{Runtime\\median (ms)} &
\shortstack{Runtime\\p95 (ms)} &
\shortstack{Ready by scheduled dispatch\\$n/N$ (\%)} &
\shortstack{Late relative to physical\\recovery loss ($n/N$)} \\
\midrule
Inspection Descent: command-only & 276 & 42.26 & 61.39 & 2760/3036 (90.91) & N/M \\
Geofence Turn: command-only & 300 & 83.70 & 104.29 & 5480/6900 (79.42) & N/M \\
\midrule
Underpass: command-only & 400 & 142.36 & 146.71 & 3600/5200 (69.23) & 0/243 \\
Underpass: RecoveryAware & 400 & 149.12 & 166.96 & 3537/5200 (68.02) & 0/243 \\
Underpass: \textbf{full} & 400 & 149.00 & 167.44 & 3547/5200 (68.21) & 0/243 \\
\bottomrule
\end{tabular}

\smallskip
\begin{minipage}{\textwidth}
\scriptsize
\textit{Definitions and scope.}
False-feasible and false-reject rates divide by physically infeasible and
feasible label counts, respectively. AE denotes absolute error.
Part A uses 19 Inspection and 20 Geofence \emph{development} scenarios,
and 400 Underpass \emph{evaluation} scenarios; Part B uses 276, 300 and
400 evaluation scenarios. Boundary errors are sampled first-loss errors,
not continuous-time deadline errors. Original eligibility conventions are
retained: Inspection excludes one initially invalid forecast; Underpass
includes one initial failure as sample zero for command/completion error.
Unobserved boundaries receive no finite error. Underpass command/completion
and recovery losses are unobserved in 82 and 157 scenarios, respectively.
N/M means not measured, not zero. Independent completion/recovery prediction
accuracy is not reported for Inspection or Geofence. Late availability
means at/after the first physically unsafe recovery sample, among the 243
Underpass scenarios with observed recovery loss. All 400 Underpass full
forecasts become ready at a tested safe recovery start.
Readiness uses the original control-step scheduling; timing was measured
offline and replayed on the simulator clock. Ready-by-dispatch counts
include cases canceled earlier. Numerical profiles coincide for
RecoveryAware and full StaleGuard, but their runtime measurements are independent.
\end{minipage}
\endgroup
\end{table*}
